\documentclass{article}
\usepackage{amsmath, amssymb}
\usepackage{bm}
\usepackage{graphicx}

\title{Low-rank Preconditioning for Systems of the Form $A + B C^{-1} B^T$}
\author{}
\date{}

\begin{document}

\maketitle

\section*{Problem Setting}

We consider the linear system:
\[
H x = (A + B C^{-1} B^T) x = b,
\]
where:
\begin{itemize}
  \item $A \in \mathbb{R}^{n \times n}$ is symmetric positive definite (SPD),
  \item $C \in \mathbb{R}^{m \times m}$ is SPD and cheap to invert,
  \item $B \in \mathbb{R}^{n \times m}$ is general,
  \item Applying $A^{-1}$ is expensive (e.g., via a Chebyshev polynomial approximation).
\end{itemize}

We aim to construct an efficient preconditioner for this system, suitable for use with the Conjugate Gradient (CG) method, while reducing the number of applications of $A^{-1}$.

\section*{Low-rank Approximation of the Schur Complement}

We define the Schur complement term:
\[
M := B C^{-1} B^T,
\quad \text{so that} \quad
H = A + M.
\]

To reduce computational cost, we approximate $M$ by a low-rank symmetric matrix:
\[
M \approx U S U^T,
\]
where $U \in \mathbb{R}^{n \times k}$ and $S \in \mathbb{R}^{k \times k}$ is symmetric positive semidefinite.

This approximation is constructed using a \textbf{randomized range finder} technique (see \cite{halko2011finding}). Given a Gaussian random matrix $\Omega \in \mathbb{R}^{n \times (k+p)}$, we compute:
\[
Y = M \Omega,
\quad \text{then perform a QR decomposition:} \quad
Y = QR,
\]
where $Q$ is an orthonormal basis of the approximate range of $M$. The projection:
\[
B_{\text{small}} = Q^T M Q
\]
is then spectrally decomposed:
\[
B_{\text{small}} = \tilde{U} S \tilde{U}^T,
\quad \text{and we set} \quad
U := Q \tilde{U}.
\]

\paragraph{Power Iteration.}
To improve the quality of the low-rank approximation—especially when the spectrum of $M$ decays slowly—we may apply a few steps of the power iteration:
\[
Y \leftarrow (M M^T)^q M \Omega,
\]
which amplifies dominant modes and improves the approximation of the leading eigenspace of $M$.

\section*{Preconditioner Using the Woodbury Identity}

Let us approximate $H \approx A + U S U^T$.
Using the Woodbury identity:
\[
(A + U S U^T)^{-1}
= A^{-1} - A^{-1} U (S^{-1} + U^T A^{-1} U)^{-1} U^T A^{-1}.
\]

This exact expression requires two applications of $A^{-1}$.
To reduce cost, we factor the formula by introducing:
\[
W := A^{-1} U,
\quad
W^T := U^T A^{-1},
\quad
T := W^T U,
\quad
M_{\text{inv}} := (S + T)^{-1}.
\]

We then write the preconditioner as:
\[
P^{-1} = A^{-1} - W M_{\text{inv}} W^T.
\]

Applying $P^{-1}$ to a vector $v$ is done as follows:
\begin{enumerate}
  \item Compute $z = A^{-1} v$,
  \item Correct: $P^{-1} v = z - (W M_{\text{inv}}W^T)v$.
\end{enumerate}

This requires only one application of $A^{-1}$ per vector, and all other operations are low-dimensional.

\section*{Remarks}

\begin{itemize}
  \item The low-rank approximation of $M$ can be significantly cheaper to construct than assembling or inverting the full matrix.
  \item The preconditioner remains symmetric positive definite if $S$ is SPD, allowing for use with CG.
  \item This approach combines randomized numerical linear algebra with classical techniques (e.g., Woodbury) to reduce iteration cost.
\end{itemize}

\section*{References}

\begin{thebibliography}{9}

\bibitem{halko2011finding}
N. Halko, P.-G. Martinsson, J. A. Tropp,\\
\textit{Finding structure with randomness: Probabilistic algorithms for constructing approximate matrix decompositions},\\
SIAM Review 53(2), 2011.

\end{thebibliography}

\end{document}
